\section{Introduction}
\label{sec:intro}

Two theories are given, each internally lawful. When may they interact, and when does
their interaction produce a new theory rather than a restatement of the two old ones?
The question arises wherever theories are combined: two axiomatizations merged into one,
two models of a system coupled through a shared instrument, two layers of description
linked by an interface. It is easy to answer carelessly. Two theories that share an
origin, a carrier or a vocabulary are said to interact; a combination that produces some
new quantity is said to be a genuine joint theory; a combination that removes an old
difficulty is said to have fewer problems. Each inference takes something real and
observable and reads more into it than follows.

This paper does two things about that situation. First, it gives a \emph{certificate
language}: typed records in which a claim about interaction must be written down, with a
separate field for each thing that is easily conflated. A record says, for instance,
whether content is expressible, present, reachable or occurrent; whether two theories are
merely in contact or have formed a composite; whether an enabling theory also determines
what it enables; whether a dependence on order has been shown or an arrow of time. A
finite reference world exhibits, for each separation, a configuration that has one
property and lacks the other. Second, it proves where such certificates have mathematical
content, by giving semantic models in which the key certificates are backed by theorems.

\paragraph{The catalog.} The certificates are organized in a \emph{catalog}: a list of 36
entries, each a claim about lawful interaction together with the certificate that must
back it, the finite evidence attached to it and its counterpart in the Lean development.
The entries fall into four clusters (access and admission, contact and join, enablement
and descent, dynamics and residue) and are anchored by one finite reference world
(\Cref{sec:ttw}). It is worth being plain about what the catalog contains. Most entries
fix what a certificate must record and when it counts as valid; they are definitions, and
the statements drawn from them are short consequences of those definitions or refusals of
credit that a definition withholds. \Cref{tab:status} classifies every result by what it
is. The mathematics lies in the semantic results of
\Cref{sec:strict-join,sec:admission}, which the introduction leads with.

\subsection{A worked example}
\label{sec:intro-example}

\paragraph{Two theories and a join.}
Let theory~$A$ observe one bit $a$ and theory~$B$ observe one bit $b$. A candidate joint
theory has states $j=(a,b,h)$ with a third bit~$h$, and maps $p_A(j)=a$ and $p_B(j)=b$
back to the two parents. Which observables of the joint theory are genuinely new?

The observable $f(j)=a$ is not new: it factors through $p_A$ on its image. The observable $f(j)=h$ is.
The two states $(0,0,0)$ and $(0,0,1)$ have the same image under $p_A$, under $p_B$, and
under the pair $(p_A,p_B)$, yet they differ in~$h$. No function of $a$, of $b$, or of
$(a,b)$ can therefore return~$h$. We call such a pair of states a \emph{split pair}.
\Cref{thm:split-criterion} shows that split pairs are exactly what stands in the way of
factorization on the reached image, so that $h$ is \emph{strict} in the sense of \Cref{def:strict-observable}
(\Cref{ex:cube}).

The pairing matters. An observable such as $a\oplus b$ factors on its image through neither
parent alone, since $a$ does not determine $a\oplus b$ and neither does~$b$, but it
factors on its image through the pair $(a,b)$. It is information that the two theories already hold jointly,
and a joint theory that only adds such observables adds nothing that the two parents do
not jointly determine. This is why strictness is tested against the pairing as well as
against each parent. What the test leaves open is stated in \Cref{rem:refinement-limit}.

\paragraph{Admitting the pieces.}
Suppose now that the joint theory is assembled by admitting items one at a time, where an
item may be admitted once none of its prerequisites is still pending. Does the order of
admission matter? If admitting an item can only enable further admissions, it does not:
two admissions available at the same moment commute, every admission run terminates, and
by Newman's lemma every maximal order ends in the same state (\Cref{thm:admission}). A
single guard that lets one admission \emph{disable} another breaks this. With items $a$
and $b$, where $b$ may be admitted only while $a$ is still pending, admitting $a$ first
strands $b$, while admitting $b$ first admits both. The two orders end in different
terminal states (\Cref{ex:disabling}).

\subsection{Main results}
\label{sec:intro-results}

We state the headline results informally here; the precise statements carry all
hypotheses.

\begin{headline}[A: strictness is non-factorization]
An observable $f$ of a join factors through a map $q$ on the image of~$q$ exactly when no
split pair for $q$ and $f$ exists (\Cref{thm:split-criterion}). Consequently $f$ is strict,
factoring on its image through neither parent nor their pairing, exactly when it has a split pair
against each of the three maps (\Cref{thm:strict-split}). The third bit of the three-bit
join is strict and the first is not (\Cref{ex:cube}). Strictness rules out relabelling and
coarsening of a parent (\Cref{prop:strict-not-cheap}); \Cref{rem:refinement-limit} states its limit.
\end{headline}

\begin{headline}[B: admission order does not matter when admission only enables]
In any system where an item may be admitted once none of its prerequisites is still pending, every run
terminates, simultaneous admissions commute, and every maximal admission order from a
given state reaches the same terminal state with the same audit value
(\Cref{thm:admission}). The proof applies Newman's lemma (\Cref{thm:newman}); local
confluence of a finite system is decidable by a finite search (\Cref{prop:finite-peaks}).
Termination cannot be dropped (\Cref{ex:four-state}), and one disabling guard produces two
terminal states (\Cref{ex:disabling}).
\end{headline}

\begin{headline}[C: a soundness check for the strictness certificate]
Let $f$ be a strict observable of a semantic join, let $q$ be a declared common refinement
and $s$ a declared schedule, and suppose it is proved that $f$ factors on its image through
neither $q$ nor $s$. Then the anti-product witness whose flags are the truth values of the
corresponding semantic predicates is valid (\Cref{thm:bridge}). Strictness does not supply
the two extra hypotheses (\Cref{rem:refinement-limit}).
\end{headline}

\begin{headline}[D: what the certificate language separates]
The certificate language keeps apart seven access coordinates, contact from join,
composite formation from parent retention, enablement from descent (a recorded claim that
an enabling theory also determines what it enables), reachability from
occurrence, and order dependence from an arrow. Each separation is exhibited in the finite
reference world by a named configuration with one property and not the other
(\Cref{sec:separations}, \Cref{tab:separations}).
\end{headline}

Theorem~A is the mathematics behind the catalog's central object, the strict join.
Theorem~B is the mathematics behind its confluence entry. Theorem~C is a soundness
check between the two levels. Once its hypotheses are granted, each flag is settled in a
line by the results behind Theorem~A or by a hypothesis, so the theorem is close to a restatement. What it
checks is that the Boolean certificate asks for nothing a semantic model cannot supply:
its validity predicate is satisfied by a genuine model, with each flag the truth value of
a proved statement.

\subsection{Map of results}
\label{sec:intro-map}

\Cref{tab:map} lists the numbered results. Each is proved where it is stated.

{\footnotesize
\begin{longtable}{@{}>{\raggedright\arraybackslash}p{3.4cm}>{\raggedright\arraybackslash}p{11.1cm}@{}}
\caption{Map of results.}\label{tab:map}\\
\toprule
\textbf{Result} & \textbf{What it establishes} \\
\midrule
\endfirsthead
\toprule
\textbf{Result} & \textbf{What it establishes} \\
\midrule
\endhead
\bottomrule
\endfoot
\Cref{thm:split-criterion} & A map factors through another on its image exactly when no split pair exists; for a finite list with decidable equality, a scan decides it. \\
\Cref{thm:strict-split} & An observable of a join is strict exactly when it has split pairs against each parent and against the pairing. \\
\Cref{prop:refinement} & An observable that factors globally through the pairing factors through every common refinement. \\
\Cref{prop:strict-not-cheap} & A strict observable is neither relabelling-only nor coarsening-only. \\
\Cref{thm:bridge} & For a strict observable that does not factor through the declared refinement and is not scheduling-only, setting each flag of an anti-product witness to its semantic truth value, with a nonempty description and an audit entry, gives a valid witness; the three-bit join is an instance. \\
\Cref{thm:newman} & A terminating, locally confluent relation is confluent (classical). \\
\Cref{cor:normal-forms} & Under confluence normal forms reachable from one state coincide, and with termination each state has exactly one. \\
\Cref{prop:finite-peaks} & On a finite list with decidable equality and a decidable step relation, reachability and local confluence are decided by enumeration. \\
\Cref{thm:admission} & In every admission system, admission runs terminate and every maximal admission order from a given state reaches the same terminal state. \\
\Cref{prop:enablement-model} & An upper step can be enabled by the lower dynamics while an upper observable does not factor through a lower quotient. \\
\Cref{prop:no-free-join} & A join that requires a capability not initially available, and alone produces it, never occurs in a lawful sequence. \\
\Cref{prop:selection} & Selection by intersection creates no fact. \\
\Cref{prop:search} & A bounded scan is exactly as strong as the coverage of its list and the closure of its horizon. \\
\Cref{prop:bootstrap} & With no admitted seed, no reachable generator and no open provision, no first extension is authorized. \\
\end{longtable}
}

\subsection{Relation to earlier work}
\label{sec:intro-earlier}

The factorization criterion of Theorem~A is the classical fact that a map is constant on
the fibres of another exactly when it factors through that map on its image; this form on
the reached image recurs throughout the Six Birds series \citep{TsiokosFoundationsIV}.
Newman's lemma is classical \citep{Newman1942}, as is its standard counterexample without
termination (\Cref{ex:four-state}) and the use of critical pairs to decide local
confluence \citep{KnuthBendix1970,Huet1980}. The semantic package of
\Cref{def:semantic-join} is the theory package of Foundations~I stripped to its observable
\citep{TsiokosFoundationsI}. New here are the semantic join and its strictness
characterization (\Cref{thm:strict-split}), which gives the catalog's anti-product
condition a meaning; the admission model with \Cref{thm:admission}, the semantic content
of the catalog's confluence entry, and its disabling counterexample; the finite decision
procedure in this setting; and the soundness check of \Cref{thm:bridge}.

The certificates sharpen records of the earlier papers. The access coordinates
\exposed, \recoverable{} and \admissible{} sharpen Foundations~IV's access quotient and
Foundations~II's admissible claims, and the split into \expressible, \present,
\reachable{} and \occurrent{} is new \citep{TsiokosFoundationsIV,TsiokosFoundationsII}.
Contact surfaces and join statuses sharpen the instrument-indexed records and status
families of Foundations~III \citep{TsiokosFoundationsIII}. The strict-join certificate
sharpens strict extension as non-factorization
\citep{TsiokosFoundationsIV,TsiokosFoundationsIII,TsiokosStrictTest}. The enablement
record sharpens the distinction between enablement and within-layer causation in the
agenthood paper \citep{TsiokosAgents}. The drive certificate sharpens the protocol trap of
Foundations~I and the no-go results on arrows from route residue
\citep{TsiokosFoundationsI,TsiokosNoGo}. The bootstrap refusal sharpens the
internalization law of Foundations~V and no-free-distinction of Foundations~IV
\citep{TsiokosFoundationsV,TsiokosFoundationsIV}. \Cref{sec:related} places the
certificate language among combination methods for logical theories, interface theories
and provenance.

\subsection{Roadmap}

\Cref{sec:related} discusses related work and \Cref{sec:setting} fixes the setting and
classifies the results by status.
\Cref{sec:strict-join} proves Theorems~A and~C, and \Cref{sec:admission} proves
Theorem~B. \Cref{sec:certificates} presents the certificate language and
\Cref{sec:separations} what it separates and refuses. \Cref{sec:ttw} describes the finite
reference world, \Cref{sec:scope} states the scope of the results, and \Cref{sec:lean}
describes the Lean formalization. The companion supplement holds the source concordance,
the full field lists of the certificates, the per-entry register, the no-go register, the
finite evidence and the replay commands.
